\documentclass[a4paper,12pt]{article}
\usepackage[utf8]{inputenc}
\usepackage[T2A]{fontenc}
\usepackage[english,russian]{babel}
\usepackage{geometry}
\geometry{left=2cm,right=2cm,top=2cm,bottom=2cm}
\usepackage{amsmath}
\usepackage{multirow}
\usepackage{array}
\usepackage{float}

\begin{document}
	
	\subsubsection*{Опыт 1. Окислительные свойства хлора на примере взаимодействия с сульфидом натрия}
	
	\subsubsection*{Ход работы}
	
	\quad К 5–6 каплям хлорной воды добавлено 3–4 капли раствора сульфида натрия и 1–2 капли разбавленной соляной кислоты. Наблюдается помутнение раствора вследствие образования осадка серы.
	
	\subsubsection*{Уравнение реакции}
	\begin{align*}
		& \text{Cl}_2 + \text{Na}_2\text{S} + 2\text{HCl} \rightarrow 2\text{NaCl} + \text{S} \downarrow + 2\text{HCl} \\
		& \text{Cl}_2 + \text{S}^{2-} \rightarrow 2\text{Cl}^- + \text{S} \downarrow
	\end{align*}
	
	\quad Соляная кислота служит для создания кислой среды и подавления диспропорционирования хлора.
	
	\subsubsection*{Опыт 2. Окислительные свойства дихромата калия с иодидом калия и сульфитом натрия}
	
	\subsubsection*{Ход работы}
	
	\quad В две пробирки налито по 8–10 капель хромовой смеси \(K_2Cr_2O_7 + H_2SO_4\). В первую добавлено 3–4 капли раствора иодида калия, во вторую — столько же раствора сульфита натрия. В первой пробирке наблюдается изменение окраски на коричневую (выделение \(I_2\)), во второй — изменение цвета с оранжевого на зеленый.
	
	\subsubsection*{Уравнения реакций}
	\begin{align*}
		& \text{K}_2\text{Cr}_2\text{O}_7 + 6\text{KI} + 7\text{H}_2\text{SO}_4 \rightarrow \text{Cr}_2(\text{SO}_4)_3 + 3\text{I}_2 + 4\text{K}_2\text{SO}_4 + 7\text{H}_2\text{O} \\
		& \text{Cr}_2\text{O}_7^{2-} + 6\text{I}^- + 14\text{H}^+ \rightarrow 2\text{Cr}^{3+} + 3\text{I}_2 + 7\text{H}_2\text{O} \\
		& \text{K}_2\text{Cr}_2\text{O}_7 + 3\text{Na}_2\text{SO}_3 + 4\text{H}_2\text{SO}_4 \rightarrow \text{Cr}_2(\text{SO}_4)_3 + 3\text{Na}_2\text{SO}_4 + \text{K}_2\text{SO}_4 + 4\text{H}_2\text{O} \\
		& \text{Cr}_2\text{O}_7^{2-} + 3\text{SO}_3^{2-} + 8\text{H}^+ \rightarrow 2\text{Cr}^{3+} + 3\text{SO}_4^{2-} + 4\text{H}_2\text{O}
	\end{align*}
	
	\subsubsection*{Опыт 3. Влияние среды на окисление сульфита натрия перманганатом калия}
	
	\subsubsection*{Ход работы}
	
	\quad В три пробирки налито по 3–4 капли раствора \(KMnO_4\). В первую добавлено 2 капли \(H_2SO_4\) (кислая среда), во вторую — 2 капли \(H_2O\) (нейтральная), в третью — 2 капли \(NaOH\) (щелочная). Во все пробирки внесен микрошпатель кристаллического \(Na_2SO_3\). Наблюдаемые изменения:
	
	\begin{itemize}
		\item Кислая среда: обесцвечивание, выделение бесцветного газа (\(SO_2\)).
		\item Нейтральная среда: бурый осадок \(MnO_2\).
		\item Щелочная среда: зеленое окрашивание (\(MnO_4^{2-}\)).
	\end{itemize}
	
	\subsubsection*{Уравнения реакций}
	\begin{align*}
		& \text{KMnO}_4 + 5\text{Na}_2\text{SO}_3 + 3\text{H}_2\text{SO}_4 \rightarrow \text{MnSO}_4 + 5\text{Na}_2\text{SO}_4 + \text{K}_2\text{SO}_4 + 3\text{H}_2\text{O} \\
		& 2\text{KMnO}_4 + 3\text{Na}_2\text{SO}_3 + \text{H}_2\text{O} \rightarrow 2\text{MnO}_2 \downarrow + 3\text{Na}_2\text{SO}_4 + 2\text{KOH} \\
		& 2\text{KMnO}_4 + \text{Na}_2\text{SO}_3 + 2\text{NaOH} \rightarrow 2\text{KNaMnO}_4 + \text{Na}_2\text{SO}_4 + \text{H}_2\text{O}
	\end{align*}
	
	\subsubsection*{Опыт 4. Восстановительные свойства металлов на примере взаимодействия цинка с сульфатом меди(II)}
	
	\subsubsection*{Ход работы}
	В пробирку с 10–15 каплями \(CuSO_4\) опущен кусочек цинка. Наблюдается выделение красного осадка меди и посветление раствора.
	
	\subsubsection*{Уравнение реакции}
	\[
	\text{Zn} + \text{CuSO}_4 \rightarrow \text{ZnSO}_4 + \text{Cu} \downarrow
	\]
	\[
	\text{Zn} + \text{Cu}^{2+} \rightarrow \text{Zn}^{2+} + \text{Cu} \downarrow
	\]
	
	\subsubsection*{Опыт 5. Восстановительные свойства соединений с элементами в низших степенях окисления}
	
	\subsubsection*{Ход работы}
	\begin{enumerate}
		\item К 3–4 каплям бромной воды добавлено 2–3 капли конц. \(NH_3\). Наблюдается обесцвечивание.
		\item К 2–3 каплям \(Na_2S\) добавлено 3–5 капель разб. \(HNO_3\). Наблюдается выделение газа и помутнение.
	\end{enumerate}
	
	\subsubsection*{Уравнения реакций}
	\begin{align*}
		& 3\text{Br}_2 + 8\text{NH}_3 \rightarrow 6\text{NH}_4\text{Br} + \text{N}_2 \\
		& 3\text{Na}_2\text{S} + 8\text{HNO}_3 \rightarrow 3\text{S} \downarrow + 2\text{NO} \uparrow + 6\text{NaNO}_3 + 4\text{H}_2\text{O}
	\end{align*}
	
	\subsubsection*{Общий вывод}
	
	\quad В ходе лабораторной работы №7 были изучены окислительно-восстановительные свойства различных веществ. Установлено, что:
	
	\begin{itemize}
		\item Хлор проявляет сильные окислительные свойства, окисляя сульфид-ионы до элементарной серы.
		\item Дихромат калия в кислой среде является окислителем для иодид- и сульфит-ионов.
		\item Продукты восстановления перманганата калия зависят от среды: в кислой среде образуется \(Mn^{2+}\), в нейтральной — \(MnO_2\), в щелочной — \(MnO_4^{2-}\).
		\item Цинк как активный металл вытесняет медь из раствора её соли.
		\item Аммиак и сульфид натрия проявляют восстановительные свойства по отношению к брому и азотной кислоте соответственно.
	\end{itemize}
	
	\quad Экспериментально подтверждены теоретические положения об окислительно-восстановительных процессах, влиянии среды на их протекание и свойствах важнейших окислителей и восстановителей.
	
\end{document}